\documentclass[12pt]{article}
\usepackage[utf8]{inputenc}
\usepackage[T1]{fontenc}
\usepackage{hyperref}
\title{Zokhuo Language Resources for K-12 Curriculum Development}
\author{Compiled by Claude Sonnet 5.5 (claude-sonnet-5-5) for Andr\'as Kornai}
\date{}
\begin{document}
\maketitle

\section{Introduction}
Zokhuo (ISO 639-3 yzk; Glottocode zokh1238) is a Loloish (Sino-Tibetan) language of Yunnan, China. Sources place speakers in southeast Wenshan County and south Yanshan County, with about 13,000 speakers reported for 2011. Speakers are officially classified as Yi. Alternate names include Cowtail Phula, Niuweiba Phula, Nimitso and Zuoke. Reports differ on vitality, from threatened to thriving. My searches found no school materials in Zokhuo, and no primer, textbook or dictionary.

\section{Primary Educational Resources}
No primers, textbooks or school programmes in Zokhuo turned up in my searches. Any curriculum would have to be built from scratch, with community involvement.

\section{Linguistic and Reference Materials}
\begin{itemize}
\item Wikipedia, Zokhuo language (overview, classification, location, speaker numbers): \url{https://en.wikipedia.org/wiki/Zokhuo_language}
\item Wikitongues, Zokhuo (language profile): \url{https://wikitongues.org/languages/yzk/}
\end{itemize}

\section{Religious and Community Resources}
\begin{itemize}
\item Joshua Project, Zokhuo in China (people-group profile): \url{https://joshuaproject.net/people_groups/20617}
\end{itemize}

\section{Digital and Audio Resources}
\begin{itemize}
\item Global Recordings Network, Zokhuo language page (the search result gave only the page, and I did not check which audio it holds): \url{https://globalrecordings.net/en/language/yzk}
\end{itemize}

\section{Curriculum Development Context}
Little published material exists for Zokhuo. The language is reported as the first language of all adults, but it is not known to be taught in schools. A K-12 curriculum would need new documentation work: an orthography, wordlists and recorded texts. It would also need to account for the language's close relation to Khlula and for the Chinese-medium school system.

\end{document}